\subsection{模；向量空间；线性组合}

定义 1. 给定一个环 A，A 上的左模（或左 A-模）是指一个集合 E，连同通过给出以下内容而定义的代数结构：

\begin{enumerate}
\def\labelenumi{(\arabic{enumi})}
\item
  E 上的一个交换群律（以下均按加法记法书写）；
\item
  一个作用律 \((\alpha, x) \mapsto \alpha \tau x\) ，其算子域为环 A，且满足下列公理：
\end{enumerate}

\((\mathbf{M}_{\mathbf{I}})\alpha \tau (x + y) = (\alpha \tau x) + (\alpha \tau y)\) 对所有 \(\alpha \in \mathbf{A}, x \in \mathbf{E}, y \in \mathbf{E}\) ；

\((M_{\mathrm{II}}) (\alpha + \beta) \tau x = (\alpha \tau x) + (\beta \tau x) \text{ 对所有 } \alpha \in A, \beta \in A, x \in E;\)

\((M_{\mathrm{III}})\alpha \tau (\beta \tau x) = (\alpha \beta)\tau x\) 对所有 \(\alpha \in A,\beta \in A,x\in E;\)

\((M_{\mathrm{IV}}) \ 1 \tau x = x \ for \ all \ x \in E.\)

公理（M \(_{I}\) ）意味着左 A-模 E 的外合成律关于 E 上的加法是分配的；因此，模是一个带算子的交换群。

如果在定义1中，公理 \((\mathbf{M}_{\mathrm{III}})\) 被替换为

\[
\left(\mathrm{M} _ {\mathrm{III}} ^ {\prime}\right) \alpha \tau (\beta \tau x) = (\beta \alpha) \tau x for all \alpha \in \mathrm{A}, \beta \in \mathrm{A}, x \in \mathrm{E},
\]

E 连同如此定义的代数结构是 A 上的右模或右 A-模。

当谈到 A-模（左模或右模）时，环 A 的元素常称为标量。

通常，左 A-模（相应地，右 A-模）的外合成律按乘法记法书写，算子写在左边（相应地，右边）；此时条件 \((\mathbf{M}_{\mathrm{III}})\) 写作 \(\alpha(\beta x) = (\alpha\beta)x\) ，条件 \((\mathbf{M}_{\mathrm{III}}^{\prime})\) 写作 \((x\beta)\alpha = x(\beta\alpha)\) 。

若 \(A^{0}\) 表示 A 的反环（I，§8，第 3 号），则环 A 上的每个右模 E 都是环 \(A^{0}\) 上的左模。由此可知，模的性质的阐述可以系统地只限于左模或只限于右模来进行；在 §§1 与 §§2 中，我们一般对左模给出这种阐述，而当我们说到一个模（不指明是哪一类）时，指的是外合成律按乘法记法书写的左模。当环 A 交换时，关于 A 的右模与左模的概念是相同的。

对于所有 \(\alpha \in \mathbf{A}\) ，映射 \(x \mapsto \alpha x\) 将 A-模 \(\mathbf{E}\) 映入自身，称为以 \(\alpha\) 为位似比的 \(\mathbf{E}\) 的位似（I，§4，第 2 号）；由 \((\mathbf{M}_{\mathbf{I}})\) ，位似是 \(\mathbf{E}\) 的交换群结构（不带算子）的自同态，但一般不是 \(\mathbf{E}\) 上模结构的自同态（I，§4，第 2 号；参见 II，§1，第 2 号和第 13 号）。因此 \(\alpha 0 = 0\) 且 \(\alpha(-x) = -(\alpha x)\) ；若 \(\alpha\) 是 \(\mathbf{A}\) 的一个可逆元素，则位似 \(x \mapsto \alpha x\) 是 \(\mathbf{E}\) 的交换群结构（不带算子）的自同构 ，因为由关系 \(y = \alpha x\) 并根据 \((\mathbf{M}_{\mathbf{IV}})\) 可得 \(x = \alpha^{-1}(\alpha x) = \alpha^{-1}y\) 。

类似地，根据 \((\mathbf{M}_{\mathrm{II}})\) ，对所有 \(x \in E\) ，映射 \(\alpha \mapsto \alpha x\) 是加法群 A 到交换群（不带算子）E 的同态；因此 0x = 0 且 \((- \alpha)x = -(\alpha x)\) ；此外，由 \((\mathbf{M}_{\mathrm{IV}})\) ，对每个整数 \(n \in Z\) ，有 \(n.x = (n.1)x\) 。

当环 A 仅由元素 0 组成时，每个 A-模 E 也仅由元素 0 组成，因为此时 A 中 \(1 = 0\) ，从而对所有 \(x \in \mathbf{E}\) ，有 \(x = 1.x = 0.x = 0\) 。

例。(1) 设 \(\phi\) 是环 A 到环 B 的一个同态；映射 \((a, x) \mapsto \phi(a)x\) （相应地 \((a, x) \mapsto x\phi(a)\) ）将 A × B 映入 B，在 B 上定义了一个左（相应地，右）A-模结构。特别地，若取 \(\phi\) 为 A 上的恒等映射，则在 A 上得到一个典范的左（相应地，右）A-模结构；为避免混淆，具有此结构的集合 A 记作 A \(_s\) （相应地，A \(_d\) ）。

\begin{enumerate}
\def\labelenumi{(\arabic{enumi})}
\setcounter{enumi}{1}
\item
  在交换群 G（按加法记法书写）上，由外合成律 \((n, x) \mapsto n.x\) （I，§3，第 1 号）定义的带算子群结构是有理整数环 Z 上的模结构。
\item
  设 E 为按加法记法书写的交换群， \(\mathcal{E}\) 为 E 的自同态环（I，§8，第 3 号：回顾，两个自同态的乘积 fg 按定义是复合自同态 \(f \circ g\) ）。外合成律 \((f, x) \mapsto f(x)\) （介于算子 \(f \in \mathcal{E}\) 与元素 \(x \in E\) 之间）在 E 上定义了一个典范的左 \(\mathcal{E}\)-模结构。
\end{enumerate}

现在考虑一个环 A，并假设在 E 上给定了一个左（相应地，右）A-模结构；对所有 \(\alpha \in \mathbf{A}\) ，位似 \(h_{\alpha}: x \mapsto \alpha x\) （相应地 \(x \mapsto x\alpha\) ）属于 \(\mathcal{E}\) ；映射 \(\phi: \alpha \mapsto h_{\alpha}\) 是环 A（相应地，反环 \(\mathbf{A}^0\) ）到环 \(\mathcal{E}\) 的同态，且根据定义 \(\alpha x = (\phi(\alpha))(x)\) （相应地 \(x\alpha = (\phi(\alpha))(x)\) ）。反之，给定环同态 \(\phi: \mathbf{A} \to \mathcal{E}\) （相应地 \(\phi: \mathbf{A}^0 \to \mathcal{E}\) ），由上述公式便在 E 上定义了一个左（相应地，右）A-模结构。换言之，在加法群 E 上（其加法律为给定的群律）给定一个左（相应地，右）A-模结构，等价于给定一个环同态 \(\mathbf{A} \to \mathcal{E}\) （相应地 \(\mathbf{A}^0 \to \mathcal{E}\) ）。

定义 2. 域 K 上的左（相应地，右）向量空间是指一个左（相应地，右）K-模。

向量空间的元素有时称为向量。

例。(4) 一个域关于其任意子域既是左向量空间又是右向量空间。

*(5) 三维实数空间 \(\mathbf{R}^3\) 是关于实数域 \(\mathbf{R}\) 的向量空间，乘积 \(tx\) （ \(t\) 为实数， \(x\) 为坐标 \(x_1, x_2, x_3\) 的点）是坐标为 \(tx_1, tx_2, tx_3\) 的点。类似地，定义在任意集合 \(\mathbf{F}\) 上的实值函数的集合是关于 \(\mathbf{R}\) 的向量空间，乘积 \(tf\) （ \(t\) 为实数， \(f\) 为函数）是实值函数 \(x \mapsto tf(x)\) 。*

对 A-模 E 的元素的两个有限支集族 \((x_{i})_{i\in I}, (y_{i})_{i\in I}\) （I，§2，第 1 号），下列等式成立：

\begin{enumerate}
\def\labelenumi{(\arabic{enumi})}
\tightlist
\item
\end{enumerate}

\[
\sum_ {i \in I} (x _ {i} + y _ {i}) = \sum_ {i \in I} x _ {i} + \sum_ {i \in I} y _ {i}\tag{2}
\]

\[
\alpha . \sum_ {i \in I} x _ {i} = \sum_ {i \in I} (\alpha x _ {i}) \quad \text { for   all } \alpha \in A;
\]

通过取 I 的一个包含 \((x_{i})\) 与 \((y_{i})\) 的支集的有限子集 H，这些等式立即化为关于有限和的相应等式。

定义 3. 称元素 \(x\) （A-模 \(\mathbf{E}\) 的元素）为系数在 \(\mathbf{A}\) 中的、元素族 \((a_{i})_{i\in I}\) （其元素属于 \(\mathbf{E}\) ）的线性组合，如果存在一个族 \((\lambda_{i})_{i\in I}\) ，其元素属于 \(\mathbf{A}\) ，具有有限支集，使得 \(x = \sum_{i\in I}\lambda_{i}a_{i}\) 。

一般地，满足此条件的族 \((\lambda_{i})\) 可以有若干个不同的（参见第 11 号）。

注意，0 是 E 的元素的空族的唯一线性组合（按 I，§2，第 1 号的约定）。

